%% kapa: exam=1 round=24 year=2013 period=2 form=A subject=accounting subject_ko=회계학 seq=1 number=1
%% source: Q-Net 감정평가사 기출문제 (cst00309) artlSeq=5073802 file=extracted/1008/5073802/2013년 제24회 감정평가사 1차시험 기출문제/2교시(A형).hwp
%% conversion: kordoc 4.24.51 -> tools/gen1.py (auto); review=proofread-sonnet; flags=none
\begin{question}{1}
재무보고를 위한 개념체계 상 재무제표 구성요소 및 인식에 관한 설명으로 옳지 않은 것은?
\begin{choices}
  \item 자산은 과거 사건의 결과로 기업이 통제하고 있고 미래경제적효익이 기업에 유입될 것으로 기대되는 자원이다.
  \item 부채는 현재 의무의 이행에 따라 경제적효익을 갖는 자원의 유출 가능성이 높고 결제될 금액에 대해 신뢰성 있게 측정할 수 있을 때 재무상태표에 인식한다.
  \item 비용은 자산의 감소나 부채의 증가와 관련하여 미래경제적효익이 감소하고 이를 신뢰성 있게 측정할 수 있을 때 포괄손익계산서에 인식한다.
  \item 수익은 자산의 증가나 부채의 감소와 관련하여 미래경제적효익이 증가하고 이를 신뢰성 있게 측정할 수 있을 때 포괄손익계산서에 인식한다.
  \item 자본은 소유주 지분의 공정가치로 측정하여 재무상태표에 인식한다.
\end{choices}
\end{question}
\answer{5}
